% Ad Familiares 9.16 --- headnote
% ad-familiares-09-16, not after mid-July 46 BC, Tusculanum

\textit{Cicero to L. Papirius Paetus, written at the Tusculan villa
not later than the middle of Quintilis (July) 46 BC --- Perseus:
\textit{in Tusculano non post med.~m.~Quint.~a.~708 (46)}. The letter
sits with 9.18 and 9.20 at the heart of the great Paetus correspondence
of mid-46 BC. Paetus has written twice in identical terms, alarmed at
something a certain Silius had said: that Cicero is being talked
about, that the wit of his after-dinner conversation is being carried
back to Caesar, that he is in some kind of danger of the new regime's
displeasure. The letter is the long, magnificent, and almost wholly
witty answer.}

\textit{It moves in three blocks. First (sections 1--6), the
political reply: Cicero has cultivated Caesar's intimates with what
\textit{art} a man can muster in this kind of weather (\textit{non
enim iam satis est consilio pugnare; artificium quoddam excogitandum
est} --- it is no longer enough to fight by strategy, some kind of
artifice has to be devised); Caesar himself is a man of acute literary
judgement who has had volumes of \textit{apophthegmata} compiled and
rejects on sight any \textit{bon mot} brought to him as Cicero's that
is not; the \textit{Oenomaus} of Accius which Paetus had quoted (some
verse on \textit{invidia}) can therefore be set aside, since the wave
that envy makes can be broken on Cicero's rock just as readily as on
the wise Greeks' who endured \textit{regna} at Athens or Syracuse.
Second (sections 7--9), the comic centrepiece: Paetus has followed
his Accius not with old-fashioned Atellan farce but, in the modern
manner, with mime, complete with a daggered \textit{popillius} and
\textit{denarius} and a dish of \textit{tyrotarichum} (cheese with
salt-fish, a poor man's casserole); Cicero parries by claiming the
new refinement of his dinner-table --- Hirtius and Dolabella are his
pupils in oratory and his masters in dining, declaiming at his house
and dining at theirs --- and threatens that when he arrives in the
country Paetus will not dare set before him \textit{a polypus to
match a red-leaded Jupiter}, the standard cheap-octopus joke
matched to the red-lead-painted statue of Capitoline Jupiter. The
appetizer course Cicero has done away with: he used to be done in
by Paetus's olives and Lucanian sausages before the main course got
its start. The closing exchange (10) dismisses the Selician villa
with the perfect Paetan epigram: \textit{salis enim satis est,
sannionum parum} --- salt there is enough, clowns too few. The
letter is the project's clearest case of Cicero deliberately
slumming into low diction with a friend who is exactly the right
audience for it; the political and the gastronomic register share
the same chiastic neatness, and the joke timing carries the
seriousness.}
